% ===== PRÉAMBULE CANONIQUE — à coller VERBATIM en tête de CHAQUE .tex =====
\documentclass[11pt,a4paper]{article}
\usepackage[utf8]{inputenc}\usepackage[T1]{fontenc}\usepackage{lmodern}
\usepackage[french]{babel}
\usepackage{amsmath,amssymb,mathtools}\usepackage{icomma}
\usepackage[version=4]{mhchem}          % \ce{...} : équations & formules
\usepackage{siunitx}\sisetup{locale=FR} % \qty{}{}, \si{}, \ang{}
\usepackage{chemfig}                    % \chemfig{...} : structures développées
\usepackage{xcolor}\usepackage{colortbl}
\usepackage{tikz}\usepackage{pgfplots}\pgfplotsset{compat=1.18}
\usepackage[most]{tcolorbox}\usepackage{enumitem}\usepackage{array,booktabs}
\usepackage[a4paper,margin=2.2cm]{geometry}
\usetikzlibrary{arrows.meta,calc,angles,quotes,positioning,patterns,backgrounds,shapes.geometric}
\definecolor{primary}{RGB}{222,49,99}\definecolor{primarydark}{RGB}{150,22,55}
\definecolor{defcol}{RGB}{0,114,178}\definecolor{methcol}{RGB}{52,92,148}
\definecolor{excol}{RGB}{0,135,90}\definecolor{attncol}{RGB}{200,40,55}
\tcbset{boxbase/.style={enhanced,breakable,boxrule=0.9pt,arc=2.5pt,
  left=3mm,right=3mm,top=2.3mm,bottom=2.3mm,fonttitle=\bfseries,coltitle=white}}
% --- boîtes TUTORIEL (noms EXACTS, ne pas renommer) ---
\newtcolorbox{cle}[1][]{boxbase,colback=defcol!6,colframe=defcol,
  title={\textsc{À retenir}\ifx\relax#1\relax\relax\else\textmd{ : \itshape #1}\fi}}
\newtcolorbox{methode}[1][]{boxbase,colback=methcol!7,colframe=methcol,sharp corners,
  title={\textsc{Méthode}\ifx\relax#1\relax\relax\else\textmd{ --- \itshape #1}\fi}}
\newtcolorbox{exemplebox}[1][]{boxbase,colback=excol!5,colframe=excol,
  title={\textsc{Exemple}\ifx\relax#1\relax\relax\else\textmd{ : \itshape #1}\fi}}
\newtcolorbox{piege}[1][]{boxbase,colback=attncol!6,colframe=attncol,
  title={\textsc{Piège}\ifx\relax#1\relax\relax\else\textmd{ : \itshape #1}\fi}}
\newtcolorbox{remarque}[1][]{boxbase,colback=black!4,colframe=black!45,
  title={\textsc{Remarque}\ifx\relax#1\relax\relax\else\textmd{ : \itshape #1}\fi}}
% --- boîtes EXERCICE / CORRIGÉ (noms EXACTS, auto-comptées) ---
\newtcolorbox[auto counter]{exo}[1][]{boxbase,colback=primary!4,colframe=primary,
  title={\textsc{Exercice}~\thetcbcounter\ifx\relax#1\relax\relax\else\textmd{ --- \itshape #1}\fi}}
\newtcolorbox[auto counter]{corrige}[1][]{boxbase,colback=excol!4,colframe=excol,
  title={\textsc{Corrigé}~\thetcbcounter\ifx\relax#1\relax\relax\else\textmd{ --- \itshape #1}\fi}}
\newcommand{\R}{\mathbb{R}}\newcommand{\N}{\mathbb{N}}\newcommand{\Z}{\mathbb{Z}}\newcommand{\ds}{\displaystyle}
% (un \usepackage{fancyhdr}+en-tête est possible ; il est ignoré côté web)
% =========================================================================
\begin{document}
\sisetup{per-mode=symbol}

\begin{corrige}[qui prédomine ?]
On compare le pH à $\mathrm{p}K_a = 4{,}8$.
\begin{enumerate}[label=\alph*)]
  \item $\mathrm{pH} = 2 < 4{,}8$ : \textbf{\ce{CH3COOH}} (l'acide) prédomine.
  \item $\mathrm{pH} = 7 > 4{,}8$ : \textbf{\ce{CH3COO-}} (la base) prédomine.
  \item $\mathrm{pH} = 4{,}8 = \mathrm{p}K_a$ : \textbf{50\,\%/50\,\%} des deux formes.
\end{enumerate}
\end{corrige}

\begin{corrige}[Henderson dans le sens direct]
$\mathrm{pH} = 4{,}8 + \log\frac{[\ce{A-}]}{[\ce{HA}]}$.
\begin{enumerate}[label=\alph*)]
  \item $\mathrm{pH} = 4{,}8 + \log 1 = 4{,}8 + 0 = \mathbf{4{,}8}$.
  \item $\mathrm{pH} = 4{,}8 + \log 10 = 4{,}8 + 1 = \mathbf{5{,}8}$.
  \item $\mathrm{pH} = 4{,}8 + \log(0{,}1) = 4{,}8 - 1 = \mathbf{3{,}8}$.
\end{enumerate}
\end{corrige}

\begin{corrige}[$\Delta\mathrm{p}K_a$ et constante d'équilibre]
\begin{enumerate}[label=\alph*)]
  \item $\Delta\mathrm{p}K_a = 10{,}4 - 2{,}1 = 8{,}3$, donc $K_{\text{équi}} = 10^{8{,}3}$ (très grand).
  \item $\Delta\mathrm{p}K_a = 4{,}8 - 6{,}4 = -1{,}6$, donc $K_{\text{équi}} = 10^{-1{,}6}$ (petit).
\end{enumerate}
\end{corrige}

\begin{corrige}[complète, équilibrée ou impossible ?]
\begin{enumerate}[label=\alph*)]
  \item $\Delta\mathrm{p}K_a = 5 > 2$ : réaction \textbf{complète}.
  \item $\Delta\mathrm{p}K_a = 1$ (entre $0$ et $2$) : réaction \textbf{équilibrée} (limitée).
  \item $\Delta\mathrm{p}K_a = -4 < -2$ : réaction \textbf{impossible} (elle irait en sens inverse).
\end{enumerate}
\end{corrige}

\begin{corrige}[réagir avec sa propre base conjuguée ?]
\begin{enumerate}[label=\alph*)]
  \item \ce{CH3COO-} capterait un \ce{H+} pour reformer \ce{CH3COOH} : l'« acide formé » serait \ce{CH3COOH}
        lui-même. Donc $\Delta\mathrm{p}K_a = 4{,}8 - 4{,}8 = 0$ et $K_{\text{équi}} = 10^0 = 1$.
  \item Non : $K_{\text{équi}} = 1$ signifie que la composition n'évolue pas. \textbf{Un acide ne réagit pas
        avec sa propre base conjuguée.}
\end{enumerate}
\end{corrige}

\end{document}
